% !TEX root = ../main.tex
\chapter{Lazy Import Surface}
\label{ch:lazy_imports}

\Lucid's top-level Python package, \code{lucid}, exposes around
three hundred names at its surface --- factories, operations,
type predicates, the \code{Tensor} class, the dtype objects, the
device objects, the autograd entry points, and a number of
subsystems re-exported as attributes (\code{lucid.nn},
\code{lucid.optim}, \code{lucid.linalg}, and others). A naive
implementation would import everything in \code{\_\_init\_\_.py},
which would force every \code{import lucid} to load the entire
engine extension, every subsystem module, and every transitively
imported file. The cold-start cost would be measured in seconds.

\Lucid\ instead uses a lazy-import surface: \code{\_\_init\_\_.py}
registers each subsystem as a deferred loader and resolves
attributes through a module-level \code{\_\_getattr\_\_} fallback.
The result is that \code{import lucid} returns in under one
hundred milliseconds on cold caches, with the heavyweight
subsystems loaded only when they are first referenced. The user
sees the same surface as before; the difference is performance.

This chapter describes the mechanism. \secref{sec:lazy_imports:why}
motivates the design. \secref{sec:lazy_imports:structure} describes
the \code{\_\_init\_\_.py} structure.
\secref{sec:lazy_imports:loaders} describes the per-subsystem
loaders. \secref{sec:lazy_imports:stubs} describes the
\code{.pyi} type stubs that preserve IDE autocomplete.
\secref{sec:lazy_imports:startup} measures the startup cost.
\secref{sec:lazy_imports:tiers} describes the Tier 1 / Tier 2 /
Tier 3 namespace layout that the lazy surface enforces.
\secref{sec:lazy_imports:adding} closes with the policy for adding
a new subsystem.

\section{Why Lazy Imports}
\label{sec:lazy_imports:why}

The cost of a non-lazy \code{import lucid} would be the sum of:

\begin{itemize}
  \item Loading the engine extension module: roughly 200--400 ms
    on cold caches (the \code{.so} file is large, the dynamic
    linker has work to do, \MLX\ initialisation runs inside).
  \item Importing every subsystem module: roughly 50--150 ms per
    subsystem, with around 15 subsystems, total in the 1--2 second
    range.
  \item Importing every model in the zoo
    (\partref{part:model_zoo}): roughly 10--30 ms per family with
    around 50 families, total in the 0.5--1.5 second range.
\end{itemize}

A pessimistic estimate puts the cold \code{import lucid} at 3--4
seconds. This is not catastrophic for a long-running training job
but it is catastrophic for several other use cases:

\begin{itemize}
  \item Interactive notebooks where the user runs many short
    cells. A 3-second import per kernel restart compounds.
  \item Command-line tools that wrap \Lucid. A model-inspection
    CLI that imports \Lucid\ at startup feels sluggish.
  \item Test suites where each test process imports \Lucid\
    fresh. A pytest run of 100 test files would add 300 seconds
    just to imports.
\end{itemize}

The lazy surface eliminates this cost for the common case where
the user does not need every subsystem. A script that uses only
\code{lucid.Tensor} and \code{lucid.nn.Linear} loads the engine
extension and the \code{nn} subsystem on first reference; the
remaining 13 subsystems and 50 model families are never loaded.

\subsection{The constraint}
\label{ssec:lazy_imports:constraint}

The lazy surface must preserve the user-facing API exactly. The
expressions \code{lucid.Tensor}, \code{lucid.zeros(3)},
\code{lucid.nn.Linear}, \code{lucid.linalg.cholesky}, and
\code{lucid.compile(model)} must all work. \code{dir(lucid)} must
return the full set of names. \code{help(lucid)} must show every
documented symbol. IDE autocomplete (Pylance, Jedi, mypy) must
work without forcing imports.

The mechanism described below satisfies all four constraints.

\section{The \_\_init\_\_.py Structure}
\label{sec:lazy_imports:structure}

\code{lucid/\_\_init\_\_.py} has three parts: the eager preamble,
the loader registry, and the deferred-resolution machinery. We
treat each in turn.

\subsection{The eager preamble}
\label{ssec:lazy_imports:eager_preamble}

A small number of objects must be available eagerly because they
are referenced during module-load time by users who unpickle
tensors, import from \code{lucid} in non-standard ways, or
otherwise interact with the module before any attribute access
triggers lazy loading:

\begin{lstlisting}[style=lucid,language=LucidPython]
# Eager: minimum surface required at import time.
from lucid._dtype  import (
    float16, float32, float64, bfloat16,
    int8, int16, int32, int64,
    bool_, complex64,
    Dtype,
)
from lucid._device import cpu, metal, Device
from lucid._tensor.tensor import Tensor
from lucid._globals import (
    get_default_dtype, set_default_dtype,
    get_default_device, set_default_device,
)

__version__ = "3.5.0"
\end{lstlisting}

This block accounts for the bulk of the 100\,ms cold-import time:
loading \code{Tensor} requires loading the engine extension, which
is the single largest contribution. Everything else in \code{lucid}
is deferred.

\subsection{The loader registry}
\label{ssec:lazy_imports:loader_registry}

The deferred symbols are registered in a dictionary that maps each
top-level name to its loader function:

\begin{lstlisting}[style=lucid,language=LucidPython]
_GROUP_LOADERS = {
    # Subsystem modules
    "nn":            "_load_nn",
    "optim":         "_load_optim",
    "autograd":      "_load_autograd",
    "linalg":        "_load_linalg",
    "fft":           "_load_fft",
    "signal":        "_load_signal",
    "special":       "_load_special",
    "distributions": "_load_distributions",
    "func":          "_load_func",
    "einops":        "_load_einops",
    "amp":           "_load_amp",
    "profiler":      "_load_profiler",
    "serialization": "_load_serialization",
    "compile":       "_load_compile",
    "models":        "_load_models",
    "utils":         "_load_utils",
    "metal":         "_load_metal",
    # Top-level functions
    "zeros":         "_load_factories",
    "ones":          "_load_factories",
    "arange":        "_load_factories",
    "tensor":        "_load_factories",
    "from_numpy":    "_load_factories",
    # ... (about 300 entries total)
}
\end{lstlisting}

The values are loader function names rather than the functions
themselves, so that the loader-function bodies (each of which
imports from a subsystem) do not run until invoked. The granularity
is per-group rather than per-symbol: \code{zeros}, \code{ones},
\code{arange}, \code{tensor}, \code{from\_numpy} all share the
\code{\_load\_factories} loader because they live in the same
subsystem and triggering one triggers all of them.

\subsection{The deferred-resolution fallback}
\label{ssec:lazy_imports:fallback}

The Python language defines that an attribute access on a module
that is not satisfied by the module's normal dictionary triggers
the module's \code{\_\_getattr\_\_} hook, if defined. \Lucid\
implements the hook to invoke the appropriate loader:

\begin{lstlisting}[style=lucid,language=LucidPython]
def __getattr__(name: str):
    loader_name = _GROUP_LOADERS.get(name)
    if loader_name is None:
        raise AttributeError(f"module 'lucid' has no attribute {name!r}")
    loader = globals()[loader_name]
    loader()  # populates globals() with the relevant symbols
    return globals()[name]
\end{lstlisting}

The pattern: on first access to \code{lucid.nn}, \code{\_\_getattr\_\_}
fires, the loader is invoked, the subsystem is imported, the
symbols are inserted into the module's globals, and the attribute
access returns successfully. Subsequent accesses to \code{lucid.nn}
hit the module dict directly (no \code{\_\_getattr\_\_} call) and
the cost is the same as for an eager import.

\subsection{The \_\_dir\_\_ extension}
\label{ssec:lazy_imports:dir}

\code{dir(lucid)} should return the full set of names, lazy or
not. The \code{\_\_dir\_\_} hook returns the union of the module's
actual globals and the keys of \code{\_GROUP\_LOADERS}:

\begin{lstlisting}[style=lucid,language=LucidPython]
def __dir__():
    return sorted(set(globals()) | set(_GROUP_LOADERS))
\end{lstlisting}

This makes tab-completion at a REPL work as expected, and makes
\code{help(lucid)} list every symbol. The cost is paid only when
the user explicitly asks for the directory listing.

\section{Subsystem Loaders}
\label{sec:lazy_imports:loaders}

Each loader function is a few lines that import from one or more
subsystem modules and insert the imports into the module globals:

\begin{lstlisting}[style=lucid,language=LucidPython]
def _load_factories():
    """Trigger import of factories subsystem and bind symbols."""
    from lucid._factories.creation import (
        zeros, ones, eye, arange, linspace, logspace, tensor,
    )
    from lucid._factories.random import (
        rand, randn, bernoulli, Generator, seed, manual_seed,
    )
    from lucid._factories.converters import (
        from_numpy, from_dlpack, to_dlpack,
    )
    for name in ("zeros", "ones", "eye", "arange", "linspace",
                 "logspace", "tensor", "rand", "randn", "bernoulli",
                 "Generator", "seed", "manual_seed",
                 "from_numpy", "from_dlpack", "to_dlpack"):
        globals()[name] = locals()[name]
\end{lstlisting}

The loader is idempotent: if invoked twice (which can happen if a
test resets the module state), the second invocation rebinds the
symbols to the same objects. Re-importing a subsystem in Python is
cheap because the module is already in \code{sys.modules}.

\subsection{Subsystem loaders as module attributes}
\label{ssec:lazy_imports:subsystem_attr}

For a subsystem module like \code{lucid.nn}, the loader does not
import individual names; it just imports the subsystem and binds
its module object as a top-level attribute:

\begin{lstlisting}[style=lucid,language=LucidPython]
def _load_nn():
    import lucid.nn
    globals()["nn"] = lucid.nn
\end{lstlisting}

The subsystem module itself can be lazy too --- \code{lucid.nn}
may not import \code{lucid.nn.functional} until referenced --- but
that is the subsystem's own concern.

\section{Type Stubs (.pyi)}
\label{sec:lazy_imports:stubs}

A lazy module is opaque to static type checkers and IDE
autocomplete by default: the type checker cannot see what symbols
\code{lucid.\_\_getattr\_\_} would return without running the
loader, and it does not run the loader. The standard solution is
the \emph{type stub}: a \code{.pyi} file that declares every
symbol's type without containing its implementation.

\code{lucid/\_\_init\_\_.pyi} contains the eager declarations:

\begin{lstlisting}[style=lucid,language=LucidPython]
from lucid._tensor.tensor import Tensor as Tensor
from lucid._dtype import (
    Dtype as Dtype,
    float16 as float16,
    float32 as float32,
    # ...
)
from lucid._device import (
    Device as Device,
    cpu as cpu,
    metal as metal,
)

# Lazy: declared here so static analysis sees them.
def zeros(*shape: int, dtype: Dtype = ..., device: Device = ...) -> Tensor: ...
def ones(*shape: int, dtype: Dtype = ..., device: Device = ...) -> Tensor: ...
def arange(start: int | float, stop: int | float | None = ...,
           step: int | float = 1, *, dtype: Dtype = ..., device: Device = ...
          ) -> Tensor: ...
# ... and so on for every name in _GROUP_LOADERS.
\end{lstlisting}

The stub is generated from the actual source code by a build-time
script and is shipped in the wheel alongside the Python sources. A
Pylance- or mypy-enabled editor sees the stub before any
\code{lucid} code runs, so autocomplete and type checking work
immediately.

\subsection{The H9 rule}
\label{ssec:lazy_imports:h9}

The stub generator must produce \emph{full} signatures, not
\code{*args, **kwargs}. Hard Rule H9 (\appref{app:hard_rules})
prohibits \code{*args/**kwargs} in stubs except for genuinely
variadic APIs. The reasoning is that a stub with \code{*args,
**kwargs} provides no static type information, defeating the
purpose of the stub.

The rule is enforced by a check in
\code{tools/check\_stubs.py} that parses every \code{.pyi} file
and flags any \code{*args} or \code{**kwargs} not accompanied by a
specific type annotation that justifies the variadic.

\section{Startup Cost Measurement}
\label{sec:lazy_imports:startup}

We measured the cold-import cost of \Lucid\ on a representative
machine (M3~Max, macOS~26.0, Python~3.14, freshly cleared
\code{\_\_pycache\_\_}):

\begin{itemize}
  \item \code{import lucid}: 92~ms.
  \item Plus first reference to \code{lucid.Tensor}: +0~ms (eager).
  \item Plus \code{lucid.nn} (first reference): +120~ms.
  \item Plus \code{lucid.optim} (first reference): +85~ms.
  \item Plus \code{lucid.compile} (first reference): +210~ms.
  \item Plus a \code{lucid.models} model factory: +180~ms.
\end{itemize}

A typical training script that uses \code{lucid}, \code{lucid.nn},
\code{lucid.optim}, and one model factory pays roughly 480~ms of
imports total --- about half what a fully eager surface would
cost. A simple inspection script that uses only \code{lucid.Tensor}
and \code{lucid.from\_numpy} pays only the 92~ms eager cost plus
50~ms for \code{\_load\_factories} on first \code{from\_numpy}.

The eager preamble's 92~ms is dominated by the engine extension
load: loading \code{libmlx.dylib} (which \Lucid\ links against) and
running its initialiser is about 60~ms; the rest is dynamic-linker
overhead for the engine \code{.so} itself. There is no reasonable
way to reduce this further without restructuring \MLX\ itself.

\section{The Tier 1 / Tier 2 / Tier 3 Namespace Layout}
\label{sec:lazy_imports:tiers}

\Lucid's user-facing namespace is organised in three tiers, with
the lazy surface enforcing the boundaries:

\begin{itemize}
  \item \textbf{Tier 1.} Names directly under \code{lucid}.
    Reserved for the \refframework-compatible top-level surface:
    factories, free-function ops, dtype and device objects, the
    \code{Tensor} class. About 300 names.
  \item \textbf{Tier 2.} Names under a subsystem module:
    \code{lucid.nn.Linear}, \code{lucid.optim.Adam},
    \code{lucid.linalg.cholesky}. About 600 names across all
    subsystems.
  \item \textbf{Tier 3.} Internal names under underscore-prefixed
    modules: \code{lucid.\_dispatch}, \code{lucid.\_C},
    \code{lucid.\_tensor}. Not part of the public API.
\end{itemize}

The lazy surface enforces the boundaries: only names that appear
in \code{\_GROUP\_LOADERS} are reachable as \code{lucid.<name>},
which keeps the Tier 1 surface auditable. A name that the
contributor wants to expose but that is not in
\code{\_GROUP\_LOADERS} simply does not appear.

\subsection{The namespace-leak rule}
\label{ssec:lazy_imports:no_leak}

A consequence of the tier structure is Hard Rule H6 (the
namespace-leak rule, conceptually): symbols defined inside a
subsystem must not leak into Tier 1 unless explicitly exposed. For
example, \code{nn.Module}, \code{nn.Parameter}, \code{nn.Linear},
\code{optim.Adam} are subsystem types, and \Lucid\ does not expose
them at the top level.

The rule is enforced by the same check that validates the lazy
surface: if a contributor accidentally adds \code{Module} to
\code{\_GROUP\_LOADERS}, the static check rejects it. This is
mostly a documentation discipline; the practical case for not
exposing \code{nn.Module} at \code{lucid.Module} is that it is
ambiguous (which subsystem's \code{Module}?).

\section{Adding a New Subsystem}
\label{sec:lazy_imports:adding}

The protocol for exposing a new subsystem at the lazy surface:

\begin{enumerate}
  \item Implement the subsystem under \code{lucid/<subsystem>/}.
  \item Write a loader function \code{\_load\_<subsystem>} in
    \code{\_\_init\_\_.py} that imports the subsystem and binds
    it to the module globals.
  \item Add an entry to \code{\_GROUP\_LOADERS} mapping the
    subsystem name to the loader.
  \item Regenerate the type stub
    (\code{tools/build\_lucid\_pyi.py}) so that the new symbols
    appear in \code{lucid/\_\_init\_\_.pyi}.
  \item Add a test that imports the new subsystem lazily and
    verifies the imports complete.
\end{enumerate}

The protocol has been invoked five times in \Lucid's history: once
each for \code{einops}, \code{amp}, \code{distributions},
\code{compile}, and \code{models}. Each invocation took roughly a
half-day of work concentrated in steps 4 and 5; steps 1--3 are
mechanical.

\section{Summary and Forward References}
\label{sec:lazy_imports:summary}

The lazy import surface in \code{lucid/\_\_init\_\_.py} reduces
cold \code{import lucid} from a worst-case 3--4 seconds to a
measured 92~ms, by deferring every subsystem import until first
reference. The mechanism is a \code{\_GROUP\_LOADERS} dictionary
plus a module-level \code{\_\_getattr\_\_} fallback, with
\code{\_\_dir\_\_} extended so \code{dir(lucid)} returns the full
surface and \code{.pyi} stubs preserving IDE autocomplete. The
Tier 1 / Tier 2 / Tier 3 namespace structure is enforced by the
same mechanism.

This chapter closes \partref{part:system}. The system architecture
chapters (the layered DAG, the tensor model, the allocator, the
dispatcher, the bridge boundaries, the pybind11 layer, and the
lazy surface) together describe the engine's static structure.
\partref{part:ops_kernels} now describes the operations themselves:
the op registry, the kernel abstraction, and the per-category
operation sets.

\begin{lucidnote}
For the user: the lazy surface is invisible. \code{import lucid}
just works; \code{lucid.nn.Linear} just works. The user does not
notice the deferred loading except that \code{import lucid} is
fast. For the contributor, the lazy surface is the disciplinary
mechanism that keeps the Tier 1 namespace from accumulating
inadvertent symbols.
\end{lucidnote}
